\documentclass[10pt,a4paper]{amsart}
\usepackage[margin=19mm]{geometry}
\usepackage{amsmath,amssymb,mathtools}
\usepackage[scr=rsfs,cal=euler]{mathalfa}
\usepackage{xfrac}
\usepackage[unicode,hidelinks,bookmarks=false]{hyperref}
\newcommand{\ie}{i.e.\ }
\newcommand{\cf}{cf.\ }
\newcommand{\resp}{respectively\ }
\newcommand{\Subsection}{\subsection}
\providecommand{\nobreakdash}{\nobreak\hbox{-}}
\setlength{\emergencystretch}{2em}
\multlinegap0pt
\usepackage{amssymb}
\usepackage{mathtools}
\usepackage[shortlabels]{enumitem}
\usepackage{tikz}
\usepackage{algorithm}\usepackage{algpseudocode}
\newtheorem{lemma}{Lemma}
\newtheorem{theorem}{Theorem}
\newcommand{\bin}[1]{\binom{#1}{2}}
\title{On the Complexity of Minimum Riesz s-Energy Subset Selection in Euclidean and Ultrametric Spaces}
\author{Michael T. M. Emmerich, Ksenia Pereverdieva, Andr\'e Deutz}
\date{Source: arXiv:2605.04715v1 (CC BY 4.0)}
\begin{document}
\maketitle

\noindent\emph{Source and licence.} On the Complexity of Minimum Riesz s-Energy Subset Selection in Euclidean and Ultrametric Spaces. Version arXiv:2605.04715v1 (CC BY 4.0), distributed by arXiv under the Creative Commons Attribution 4.0 International licence, http://creativecommons.org/licenses/by/4.0/, as recorded in the arXiv licence metadata for this exact version and shown on the abstract page https://arxiv.org/abs/2605.04715v1. Full licence notice: \texttt{licenses/arxiv-2605.04715-CC-BY-4.0.txt}.\par\smallskip

\noindent\textit{Editorial scope.} Counted unit: the article's theorem thm:gis (source0.tex lines 452--454). The article's complete original proof of it is delivered with it (source0.tex lines 456--497, one \texttt{proof} environment of 42 lines). No mathematics is written, repaired or re-derived: the statement and the proof are the article's own lines, in the article's own notation, and no retained line is substituted or re-flowed.  Retained uncounted as the source-local context the proof needs: lines 379--388; lines 413--427.  Nothing was omitted from any retained span, so the package carries no omission record.  No retained line contains a citation, so the wrapper carries no bibliography and no bibliography entry.  No retained line contains a figure environment, an image-inclusion command or drawing code, so no image is delivered and the package declares no asset dependency.  Editorial formatting only, in addition: an AMS \texttt{amsart} preamble that restates the article's own declarations and macros used by the retained lines, namely the environments \texttt{lemma}, \texttt{theorem} on separate counters, together with the standard packages those lines need.  The retained environments print in this document as Lemma 1, Theorem 1 in order of appearance, whereas the article prints the same items with the numbering of its own full document.  The complete native source of the article as distributed in the arXiv e-print for this exact version is bundled unchanged as \texttt{sources/199888/source0.tex}.
\begin{lemma}\label{lem:dmin-dmax-monotone}
Let $\emptyset\neq S\subseteq P$. Then
\[
\delta_{\max}(S)\le \delta_{\max}(P)<\delta\le D_{\min}(P)\le D_{\min}(S).
\]
In particular, if both forbidden and admissible pairs occur in $P$, then
\[
\delta_{\max}<D_{\min}.
\]
\end{lemma}

\begin{lemma}[Auxiliary exponent choice]\label{lem:aux-exp-choice}
Let $c>0$ and let $0<a<b$. Then the inequality
\[
cb^{-u}<a^{-u}
\]
holds precisely for
\[
u>\frac{\log(c)}{\log(b)-\log(a)}.
\]
In particular, if $u_0$ is any real number with
\[
u_0>\frac{\log(c)}{\log(b)-\log(a)},
\]
then $cb^{-u}<a^{-u}$ for every $u\ge u_0$.
\end{lemma}

\begin{theorem}\label{thm:gis}
RSSP-Decision in the Euclidean plane is NP-hard when $s$ is part of the input.
\end{theorem}

\begin{proof}
Let $(P,\delta,k)$ be an instance of geometric independent set. If either
$(P\times P)_{<\delta}=\emptyset$ or $(P\times P)_{\ge\delta}=\emptyset$, then
the instance is trivial and can be handled directly. We therefore assume that
both forbidden and admissible pairs occur. By
Lemma~\ref{lem:dmin-dmax-monotone}, we have $\delta_{\max}<D_{\min}$.

By comparing squared distances, compute in polynomial time the realizing pairs for $\delta_{\max}$ and $D_{\min}$, and represent these distances exactly as described above. Choose the positive integer
\[
s:=1+\left\lceil\frac{\log \bin{k}}{\log(D_{\min}/\delta_{\max})}\right\rceil,
\qquad
T:=\bin{k}D_{\min}^{-s}.
\]
The exponent $s$ is encoded in binary and $T$ is encoded by the displayed expression. By Lemma~\ref{lem:aux-exp-choice} with $c=\bin{k}$, $a=\delta_{\max}$, and $b=D_{\min}$, this choice guarantees
\[
\bin{k}D_{\min}^{-s}<\delta_{\max}^{-s}.
\]
We claim that for every $k$-subset $S\subseteq P$:
\begin{enumerate}
\item if $S$ is $\delta$-independent, then $E_s(S)\le T$;
\item if $S$ is not $\delta$-independent, then $E_s(S)>T$.
\end{enumerate}

If $S$ is $\delta$-independent, then every selected pair lies in
$(S\times S)_{\ge\delta}$, so every selected distance is at least
$D_{\min}(S)\ge D_{\min}$ by Lemma~\ref{lem:dmin-dmax-monotone}. Hence
\[
E_s(S)\le \bin{k}D_{\min}^{-s}=T.
\]
Conversely, if $S$ is not $\delta$-independent, then some selected pair has
distance strictly smaller than $\delta$, and therefore at most
$\delta_{\max}(S)\le \delta_{\max}$ by
Lemma~\ref{lem:dmin-dmax-monotone}. Hence
\[
E_s(S)\ge \delta_{\max}^{-s}.
\]
Therefore $T<\delta_{\max}^{-s}$. Hence every non-independent $k$-subset has
energy strictly larger than $T$.

Thus the GIS instance is a YES-instance if and only if the constructed RSSP
instance is a YES-instance.
\end{proof}

\end{document}
